\section{PSPACE-complete}
The preceding section established that Sokoban is an $NP\text{-hard}$ problem, following the proof of Fryers et al.\ \cite{fryers1995sokoban, cook2023complexity}. In addition to being an NP-hard game, Sokoban was also found to be PSPACE-complete \cite{culberson1997sokoban}. A PSPACE problem is a decision problem that can be solved by a Turing machine using a polynomial amount of memory space. If every problem in PSPACE can be reduced to a given problem $X$, then $X$ is PSPACE-hard. A problem is PSPACE-complete if and only if it is in PSPACE and is PSPACE-hard.

Culberson demonstrated that Sokoban is a PSPACE-complete problem by recreating a finite-tape Turing machine in Sokoban \cite{culberson1997sokoban}. To create the emulator, the researcher relied on two basic properties of the game. The first is the notion of unrecoverable configuration. This is when it is not possible to move a box to its goal position, either because the box was placed in a dead-end location or because the goal position is impossible to reach due to its placement. The second is the fact that the player's movement is also bound by the geometry of the puzzle. Several gadgets are then introduced, similarly to the NP-hard proof. Each gadget is a section of a Sokoban puzzle. These gadgets are combined into Cellular Units, each representing one cell of the finite Turing machine's tape. Culberson proves that the puzzle created from these Cellular Units has a solution if and only if the emulated Turing machine halts in an ``accept'' state.